% =============================================================
%  ch18_projekt4.tex  --  Projekt 4: CAD-Posenerkennung und Greifen
% =============================================================

\chapter{Projekt 4 -- CAD-basierte Posenerkennung und Greifen}

\textbf{Ziel:} Bauteile per Stereokamera über sich matchende CAD-Modelle
erkennen, deren 6D-Lage bestimmen und das Greifen daran anpassen.

\section{ROS\,2-Bausteine}
\begin{longtable}{@{}p{2.6cm}p{12cm}@{}}
\toprule
\textbf{Typ} & \textbf{Konkret} \\
\midrule
Nodes       & \code{stereo\_camera} $\cdot$
              \code{part\_detector} (ORB/SIFT-Matching gegen gerenderte
              CAD-Ansichten \emph{oder} Punktwolke+ICP) $\cdot$
              \code{pose\_estimator} (\code{solvePnP}/ICP $\rightarrow$ 6D-Pose) $\cdot$
              \code{moveit2} (Greifplanung) $\cdot$
              \code{robot\_controller} \\
Topics/Msgs & \code{vision\_msgs/Detection3DArray} $\cdot$
              Pose über \code{/tf} \\
Actions     & MoveIt\,2 \code{MoveGroup} (Planen + Ausführen) \\
Parameter   & CAD-Modellpfade, Matching-Schwellen, ICP-Iterationen,
              Greifvorgaben \\
\bottomrule
\end{longtable}

\begin{center}
\begin{tikzpicture}[node distance=5mm and 12mm]
\node[node] (cam) {/stereo\_camera};
\node[node,right=of cam] (det) {/part\_detector\\\scriptsize CAD-Matching};
\node[node,right=of det] (pe)  {/pose\_estimator\\\scriptsize PnP / ICP};
\node[node,right=of pe]  (mv)  {/moveit2};
\node[node,below=8mm of mv] (rc) {/robot\_controller};
\draw[flow] (cam) -- (det);
\draw[flow] (det) -- (pe);
\draw[flow] (pe)  -- node[topic,fill=white,inner sep=1pt]{\scriptsize 6D pose (tf)} (mv);
\draw[flow] (mv)  -- (rc);
\end{tikzpicture}
\captionof{figure}{Projekt~4: CAD-Matching $\rightarrow$ 6D-Posenschätzung
$\rightarrow$ MoveIt\,2-Greifplanung, angepasst an die erkannte Werkstücklage.}
\end{center}

\section{Zwei Wege zur 6D-Pose}
Für das CAD-Matching gibt es einen klassischen und einen lernbasierten Pfad --
beide enden im selben \code{pose\_estimator} und liefern $(\bm{R},\bm{t})$ über
\code{/tf} an MoveIt\,2:
\begin{enumerate}
\item \textbf{Klassisch (kein Training):} ORB/SIFT-Matching gegen gerenderte
  CAD-Ansichten $\rightarrow$ \code{solvePnPRansac}, \emph{oder} Stereo-Punktwolke
  $+$ ICP gegen die CAD-Punktwolke. Schnell aufzusetzen, aber fragil bei
  texturlosen/glänzenden/symmetrischen Teilen.
\item \textbf{Gelernte Surface Embeddings (SurfEmb-Linie):} das im Theoriekapitel
  beschriebene Zwei-Turm-Netz. Robuster bei genau diesen Industrieteilen, dafür
  Trainingsaufwand. \textbf{Empfehlung:} klassisch starten, lernbasiert
  nachrüsten, wenn die Teile die klassischen Deskriptoren überfordern.
\end{enumerate}
Die theoretischen Grundlagen beider Wege (Rotationsrepräsentationen, kontrastives
Lernen, Two-Tower-Modell, \code{PnP}+\code{RANSAC}, grobe/feine Registrierung,
PointNet++) stehen im Theoriekapitel zur 6-DoF-Posenschätzung.

\section{Inferenz-Pipeline (lernbasiert)}
\begin{enumerate}
\item \textbf{Detektion/Segmentierung} des Objekts im Stereobild $\rightarrow$
  Bildausschnitt (RGB oder RGB-D) plus bekanntes CAD-Modell als Eingang.
\item \textbf{Embedding:} Query-Tower bettet die Pixel ein, Key-Tower die
  CAD-Oberflächenpunkte -- in denselben Merkmalsraum.
\item \textbf{Matching} im Embedding-Raum (nächster Nachbar) $\rightarrow$ dichte
  2D-3D-Korrespondenzen.
\item \textbf{Pose-Hypothesen:} \code{PnP}+\code{RANSAC} sampeln und bewerten
  (Scoring), beste Hypothese wählen.
\item \textbf{Verfeinerung} per ICP $\rightarrow$ finale $(\bm{R},\bm{t})$,
  Modell$\rightarrow$Roboterframe.
\end{enumerate}

\noindent\textbf{Punktwolken-Zusatzschritte} (3D-3D-Pfad): Hintergrund entfernen,
Voxel-Downsampling, statistischer Ausreißerfilter -- vor der Registrierung.

\begin{infobox}[title=Labels sind quasi gratis]
Das CAD-Modell wird in bekannten Posen gerendert $\rightarrow$ Ground-Truth-Pose
\emph{und} die Pixel-Punkt-Korrespondenzen liegen automatisch vor. Kein manuelles
Annotieren, keine Symmetrie-Labels. Das synthetische Training wird per
\textbf{Domain Randomization} und PBR-Rendering (Isaac Sim) gegen den
Sim-to-Real-Gap abgehärtet, optional mit kleinem realem Feintuning-Satz.
\end{infobox}

\section{Stereo als metrische Skala}
Die Stereokamera ist hier nicht nur Bildquelle: Sie löst die
\textbf{Skalenmehrdeutigkeit} der Mono-Ansicht. Tiefe geht entweder als
RGB-D-Zusatzkanal ins Netz oder wird zur Punktwolke rückprojiziert. Erst dadurch
steht $(\bm{R},\bm{t})$ in echten Millimetern -- die Voraussetzung dafür, dass
MoveIt\,2 überhaupt greifen kann. \textbf{Vorsicht:} glänzendes Metall liefert
ungültige Tiefe; saubere Kalibrierung, Rektifizierung und Tiefen-Entrauschung
sind Pflicht.

\section{Bewegte Teile: Pose-Tracking mit Kalman}
Liegt das Teil nicht still (Förderband, armmontierte Kamera), ist die bildweise
Netz-Pose verrauscht und zittrig. Ein \textbf{Kalman-Filter} (vgl.\ Projekt~1/3
und das Kalman-Kapitel) macht aus Einzelschuss ein \emph{Tracking}: Zustand =
Pose (ggf.\ + Geschwindigkeit), Messung = bildweise Netz-Pose;
Prädiktion~$\rightarrow$~Korrektur glättet Jitter, verwirft Ausreißer-Frames und
überbrückt kurze Verdeckungen. Da Rotationen nichtlinear sind: EKF/UKF oder
Filtern auf einer geeigneten Rotationsdarstellung. \emph{Dieselbe
Kalman-Kompetenz wie beim Tischtennisball -- nur auf die 6D-Pose statt die
Ballposition angewandt.}

\section{Hyperparameter und Evaluation}
\begin{longtable}{@{}p{4.4cm}p{9.8cm}@{}}
\toprule
\textbf{Stellgröße} & \textbf{Wirkung} \\
\midrule
Embedding-Dimension & Ausdrucksstärke des Merkmalsraums \\
Kontrastive Temperatur $\tau$ & Schärfe der Zuordnung im Verlust \\
\#Negative / Backbone-Tiefe & Trennschärfe bzw.\ Kapazität \\
LR / Batch / Optimierer & Trainingsdynamik \\
Augmentierungsstärke & Robustheit gegen Sim-to-Real-Gap \\
\code{RANSAC}-Iterationen + Inlier-Schwelle & Ausreißerrobustheit der Posenlösung \\
\#Pose-Hypothesen / Crop-Auflösung & Genauigkeit vs.\ Laufzeit \\
\bottomrule
\end{longtable}
\noindent\textbf{Metriken:} ADD / ADD-S (ADD-S für symmetrische Teile),
geodätischer Rotations- und euklidischer Translationsfehler, RMSE; im
Forschungskontext die BOP-Metriken (VSD, MSSD, MSPD $\rightarrow$ Average Recall)
auf Datensätzen wie T-LESS/ITODD. \textbf{Industrielle KPIs:} Taktzeit/Latenz,
Pick-Erfolgsrate, Falsch-Positiv-Rate -- das sind die Zahlen, die im
Maschinenbeladungs-Einsatz zählen.
